\documentclass[12pt]{article}

\pagestyle{empty}
\setcounter{secnumdepth}{2}
%----------------------------------------------------------------------------------------
%   Packages and configurations
%----------------------------------------------------------------------------------------
\usepackage[dutch]{babel}
\usepackage{hyperref} %Allow for references
\hypersetup{
colorlinks,
citecolor=black,
filecolor=black,
linkcolor=black,
urlcolor=black
} %Set up hyperlink colours.
\newcommand{\sectionbreak}{\clearpage} % Should start every section on its own page
\usepackage{geometry} % Required to change the page size to A4
\geometry{a4paper} % Set the page size to be A4 as opposed to the default US Letter
\usepackage{chngpage}
\usepackage{appendix}
%REMOVE IF HEADER/FOOTER BROKEN
\usepackage{fancyhdr} % Required for custom headers
\usepackage{extramarks} % Required for headers and footers
\usepackage{lastpage} % Required to determine the last page for the footer
%-----------------1
\topmargin=0cm
\oddsidemargin=0cm
\textheight=22.0cm
\textwidth=16cm
\parindent=0cm
\parskip=0.15cm
\topskip=0truecm
\raggedbottom
\abovedisplayskip=3mm
\belowdisplayskip=3mm
\abovedisplayshortskip=0mm
\belowdisplayshortskip=2mm
\normalbaselineskip=12pt
\normalbaselines
\usepackage{listings}
\usepackage[svgnames,table,xcdraw]{xcolor}
\lstset { 
language=C,
frame=single,
escapeinside={\%*}{*)}, 
breaklines=true,  
backgroundcolor=\color{black!5},
basicstyle=\footnotesize,
commentstyle=\color{mygreen},
numberstyle=\tiny\color{mygray},
rulecolor=\color{black},
keywordstyle=\color{blue},
}
\definecolor{mygreen}{rgb}{0,0.6,0}
\definecolor{mygray}{rgb}{0.5,0.5,0.5}
\definecolor{mymauve}{rgb}{0.58,0,0.82}
\usepackage{wasysym}
\usepackage{lipsum}
\pagestyle{fancy}
\lhead{Guus \textsc{Beckett} \& Jim \textsc{van Abkoude}} % Top left header
\rhead{Eindrapport} % Top center header
\lfoot{\includegraphics[height=1cm]{avans}} % Bottom left footer
\cfoot{} % Bottom center footer
\rfoot{Pagina\ \thepage} % Bottom right footer
\renewcommand\headrulewidth{0.4pt} % Size of the header rule
\renewcommand\footrulewidth{0.4pt} % Size of the footer rule
\usepackage{graphicx} % Required for including pictures

\usepackage{float} % Allows putting an [H] in \begin{figure} to specify the exact location of the figure
\usepackage{wrapfig} % Allows in-line images such as the example fish picture
\usepackage{lipsum} % Used for inserting dummy 'Lorem ipsum' text into the template
\usepackage{pdfpages}
\usepackage[font={footnotesize}]{caption}
\graphicspath{ {../Images/Logos/}{../Images/} }
\setlength\parindent{0pt} % Removes all indentation from paragraphs
%\usepackage{showframe}
\newcommand*{\SignatureAndDate}[1]{%
    \par\noindent\makebox[2.5in]{\hrulefill} \hfill\makebox[2.0in]{\hrulefill}%
    \par\noindent\makebox[2.5in][l]{#1}      \hfill\makebox[2.0in][l]{Date}%
}%Signature package
\begin{document}
\begin{titlepage}
\pagenumbering{Roman}
\newcommand{\HRule}{\rule{\linewidth}{0.5mm}} % Defines a new command for the horizontal lines, change thickness here

\center % Center everything on the page

\includegraphics[height=3cm] {avans}\\% Include a department/university logo - this will require the graphicx package
\textsc{\Large Avans Hogeschool Breda}\\[0.5cm] % Major heading such as course name
\textsc{\large Netwerktechnologie}\\[0.5cm] % Minor heading such as course title
\HRule \\[0.4cm]
{ \huge \bfseries Eindrapport}\\[0.4cm] % Title of your document
\HRule \\[1.5cm]

\begin{minipage}{0.4\textwidth}
\begin{flushleft} \large
\emph{Auteurs:}\\
Guus \textsc{Beckett} \\% Your name 
Jim \textsc{van Abkoude}
\end{flushleft}
\end{minipage}
~
\begin{minipage}{0.4\textwidth}
\begin{flushright} \large
\emph{Docenten:} \\
Andries \textsc{van Dongen} \\ % Supervisor's Name
Johan \textsc{Talboom}
\end{flushright}
\end{minipage}\\[4cm]

{\large \today}\\[3cm] % Date, change the \today to a set date if you want to be precise

Versie: 0.1.0

\vfill % Fill the rest of the page with whitespace

\end{titlepage}

\clearpage
\section*{Voorwoord}
\addcontentsline{toc}{section}{Voorwoord}
Dit eindrapport zal beschrijven bla bla bla, Zeroshell bleek kut om te PXE booten, bla bla bla.
\\
Guus Beckett \& Jim van Abkoude \\
\today \\
Breda
\newpage
\tableofcontents
\newpage
\pagenumbering{arabic}
\section{Requirements}
De requirements voor dit project waren verdeeld over de 3 apparaten die in het netwerk moeten komen.
\subsection{Gateway} % (fold)
\label{sub:gateway}
Een van de door school gestelde eisen was het plaatsten van een gateway tussen het Avans hogeschool netwerk en de TSLR van dit project, dit was waarschijnlijk bedoeld als bescherming voor het Avans netwerk tegen mogelijke fouten die de student zou kunnen maken en het omzeilen van MAC adress authenticatie. Origineel was in dit project sprake van nog een vierde apparaat in het project, namelijk de imageserver. De imageservers van Avans waren niet beschikbaar gesteld aan de studenten en dus werden de verantwoordelijkheden van de imageserver naar de gateway verplaatst.
\\De gateway moet in staat zijn als zowel een DHCP en een DNS server te fungeren waarin de DNS server een eigen TLD moet hosten voor binnenin het lokale netwerk. Ook moet gateway packetfiltering toepassen.
\\Aangezien de verantwoordelijkheden van de imageserver zijn verwerkt in de gateway vallen de requirements van de imageserver nu onder de gateway. De gateway moet daarom ook een TFTP-server, NFS-server en PXE-bootserver hebben. Het imageserver deel had ook enkele bonus requirements, dit waren het draaien van een webserver met PHP en het draaien van een MySQL server op de gateway.
% subsection gateway (end)
\subsection{Thin Server Linux Router} % (fold)
\label{sub:thin_server_linux_router}
Het hoofddeel van het project was het installeren van een TSLR (Thin Server Linux Router) en deze via PXE te laten booten. Tijdens deze PXE boot moet hij zijn bootfiles verkrijgt via TFTP, een remote filesystem gebruiken via NFS en gebruik maken van een static WAN IP.
\\De TSLR moet naast PXE ook firewall en routingfuncties hebben. De routerfuncties van de TSLR gaan over het beheren van een LAN, een DMZ en het routeren van packets tussen deze 2 netwerken.
\\Zeer veel firewall verantwoordelijkheden moeten ook door de TSLR uitgevoerd worden, zoals:
\begin{itemize}
\item Blokeren van nieuwe connecties vanuit de DMZ naar het LAN.
\item Blokkeren van nieuwe connecties van WAN naar LAN.
\item Droppen van alle ICMP packets en het toelaten van alle verbindingen van LAN naar WAN.
\item Toelaten van alle verbindingen van LAN naar DMZ.
\item Instellen van static routing.
\end{itemize}
Er waren ook een paar overige requirements zoals het systeem logs remote opslaan en een SSH shell laten draaien. Naast de overige requirements waren er ook bonus requirements, dit waren het draaien van een webserver en een infopagina met status en NIC statistieken.
% subsection thin_server_linux_router (end)
\subsection{DMZ} % (fold)
\label{sub:dmz}
Een requirement van de opdracht was dat er een DMZ zou worden opgesteld en hier een enkele (1) server in zou komen. De DMZ server moest een PHP en MYSQL webserver draaien. De DMZ server heeft ook de verantwoordelijkheid een logging server te draaien, FTP en SSH server te draaien en DNS auth te regelen.
% subsection dmz (end)
\section{Vooronderzoek}
Tijdens het vooronderzoek lag de focus op het vinden van de juiste OS voor de juiste apparaten. Dit vooronderzoek bestond uit 3 delen. De eisen duidelijk opstellen, de OS bepalen van elk van de drie apparaten en een planning opstellen.\\

De eisen die naar voren zijn gekomen tijdens dit onderzoek zijn terug te vinden in hoofdstuk~\ref{sub:thin_server_linux_router} Thin Server Linux Router op pagina~\pageref{sub:thin_server_linux_router}.\\

Het kiezen van de juiste OS voor het juiste apparaat werd gedaan door middel van het gebruik van virtuele machines, waarin de verschillende OS's die geschikt leken uitvoerig werden getest. Aangezien er tijdens het vooronderzoek geen beschikking was over de hardware waarop de OS's uiteindelijk moesten gaan draaien, waren de keuzes die tijdens deze fase van het onderzoek zijn gedaan niet gericht op de hardware waarop de OSsen moesten draaien, maar alleen op de features van de gesteste OS's. De OS's die gekozen waren na dit vooronderzoek zijn:
\begin{description}
\item[Gateway - OpenBSD] De beveiligingsgerichte focus van OpenBSD maakt het uitermate geschikt om de gateway te zijn tussen het Internet en het lokale netwerk.
\item[TSLR - Zeroshell] Het feit dat Zeroshell gemaakt is om te fungeren als routersoftware maakte het een voordehandliggende keus om te gebruiken op de TSLR.
\item[DMZ server - SuperbMiniServer] De OS die op de DMZ server zou draaien moest gemakkelijk op te zetten zijn en in te stellen zijn, vandaar de keus voor deze distro, die focust op het gemakkelijk maken van het managen van een server.
\end{description}
Ten slotte is de planning gemaakt. De planning die opgesteld is focuste vooral op de volgorde van het installeren van de OS's op de verschillende apparaten. Allereerst zou de TSLR zijn OS krijgen, waarna de Gateway ge\"installeerd worden. Als laatste zou de DMZ server zijn OS krijgen. Hierna zou de PXE boot opgezet worden en de oplevering zijn van het project.
\section{Ontwerp}
\section{Implementatie}
Tijdens het implementeren en uiwerken van het project zijn de volgende stappen doorlopen
\begin{itemize}
\item Het installeren van de router op de CF kaart
\item Het installeren van de gateway
\item Het herinstalleren van de gateway
\item het installeren van de DMZ server
\item Het herinstalleren van de DMZ server
\item Het herconfigureren van de router
\item Het herinstalleren van de router
\item Het instellen van PXE op de gateway
\item De router laten booten via PXE
\item Het installeren van webservices op de DMZ
\item het configuren van de firewall en routing
\end{itemize}
Op de TSLR werd de Linux distro Zeroshell ge\"installeerd, deze werd direct op de interne CF kaart geinstalleerd. Het instellen van Zeroshell ging zeer makkelijk zoals verwacht en de router was binnen een middag klaar voor gebruik.
\\Hierna werd OpenBSD geinstalleerd op de gateway en dat leek prima te werken, de installatie verliep snel en vlekkeloos. Na de installatie werd DNSMasq ge\"installeerd om als DHCP en DNS Server dienst te doen. 
Tijdens het testen van de gateway kwamen de projectleden snel tot de conclusie dat de gateway niet correct was ingesteld, tijdens het instellen bleek ook dat OpenBSD last had met apart beheren van de netwerk interfaces. Doordat OpenBSD last had met het apart beheren van de netwerk interfaces bleek het niet haalbaar om DHCP goed in te stellen, om dit probleem alsnog op te lossen is gekozen om alsnog te wisselen naar een Linux distro. Omdat een naar een Linux distro werd geswitched werd gekozen voor Debian sinds de projectleden hier al meer ervaring mee hebben. De overstap naar Debian was snel gemaakt doordat de configuratiebestanden wel gewoon werkte op Debian.
\\Het installeren van SMS linux op de DMZ server liep niet heel gemakkelijk sinds het geen manier bood om te installeren van USB. De installatie werd dan ook snel geschrapt en werd er voor een andere OS gekozen. Na een kort overleg tussen de projectgenoten werd besloten OpenBSD te installeren ipv SMS linux aangezien deze wel via USB ge\"installeerd kon worden en OpenBSD bekend staat om gemak van installatie en veiligheid. Het installeren van OpenBSD op de DMZ server verliep zeer gemakkelijk en het instellen verliep zonder problemen.
\\Het moment dat de router opnieuw in het systeem werd ge\"introduceert bleek dat de DHCP server niet volgens standaarden functioneerde en dat deze niet gemakkelijk vervangen kon worden. Na een middag lang de router herconfigureren werkte deze nog steeds niet naar behoren. Om de nieuwe problemen met de router op te lossen werd gekozen om een andere Linux distro te installeren, ditmaal is gekozen voor Debian. De installatie van Debian op de TSLR had enkele problemen, tijdens de installatie was sprake van een tekort aan benodigde drivers. De problemen met Debian op de TSLR zetten waren gelukkig snel verholpen en de configuratie was makkelijk gedaan sinds de oude configuratiebestanden van de vorige distro nog gewoon werkte.
\\Het instellen van PXE op de gateway was zeer snel gedaan doordat gekozen is voor het gebruik van Dnsmasq als DHCP server sinds deze een TFTP server heeft ingebakken. NFS was ook gemakkelijk geinstalleerd sinds gebruik was gemaakt van een Linux en Linux standaard ondersteuning biedt voor NFS.
\\Booten via PXE was niet zeer ingewikkeld en zonder al te veel problemen verlopen. Origineel was het remote filesystem van het kaartje af gekopieerd, helaas ging dit niet helemaal goed omdat de permissies werden veranderd en sommige files daardoor niet fatsoenlijk meer werkte. De oplossing voor het probleem met de permissies was het maken van een image direct vanaf de cf kaart. Een groot probleem was dat de netwerk interfaces tijdens boot niet hetzelfde waren als na boot en daardoor het netwerk in problemen kwam, dit probleem is slechts gedeeltelijk opgelost door een bepaalde volgorde te forceren op het apparaat. De PXE boot werkte na het oplossen van deze problemen vlekkeloos.
\\De webservices op OpenBSD installeren ging zonder enige ware problemen. Door het gebruikersvriendelijke ontwerp van OpenBSD zijn package manager was het installeren van webservices ongelooflijk simpel en gedaan met letterlijk 2 commando's. Ook is vlug een PHP pagina gemaakt om te demonstreren dat de webservices werken.
\\Firewall configuratie en het defineren van static routes was relatief moeilijk om te doen sinds de projectleden weinig verstand hadden van iptables. Na wat kort onderzoek en wat rondspelen met iptables was het gelukt om het juiste verkeer door te laten en de correcte routes te defineren.
\section{Testen}
Tijdens de uitvoering van het project is het installeren en configureren constant afgewisseld met testen. Elke keer als er een wijziging gemaakt werd in een configuratie bestand werd de werking van alle features die van het gerelateerde programma verwacht werden getest op werking. Hiermee werden fouten in deze configuratiebestanden sneller ontdekt. Dit heeft uiteindelijk veel tijd bespaard doordat er nooit een moment was waarop er meer dan \'e\'en programma schuldig kon zijn voor de fouten die er op dat moment binnen het project waren.\\
\'E\'en van de weinige problemen die zijn opgetreden tijdens het testen waren problemen met het netwerk van Avans zelf. Dit netwerk was soms totaal niet bereikbaar omdat het netwerk in het projectlokaal werd afgesloten, hierdoor is het projectteam in totaal 3 uur werktijd verloren gegaan aan het wachten op de terugkeer van het netwerk tot er weer getest kon worden.
\end{document}
